%%%%%%%%%%%%%%%%%%%%%%%%%%%%%%%%%%%%%%%%%%%%%%%%%%%%%%%%%%%%%%%%%%%%%%%%%%%%%%%
% BEHAVIORAL CHANGE JOURNEY: STAGE-DEPENDENT DYNAMICS AND INTERVENTION DESIGN
%
% Style: Ernst Fehr (Zurich Tradition)
% Target: American Economic Review / Journal of Political Economy
%
% 8D Coordinates: D1=0.9, D2=0.85, D3=0.7, D4=0.8, D5=G1, D6=1.0, D7=0.1, D8=0.95
%%%%%%%%%%%%%%%%%%%%%%%%%%%%%%%%%%%%%%%%%%%%%%%%%%%%%%%%%%%%%%%%%%%%%%%%%%%%%%%

\documentclass[12pt]{article}

\usepackage{amsmath,amssymb,amsthm}
\usepackage{graphicx}
\usepackage{booktabs}
\usepackage{natbib}
\usepackage{setspace}
\usepackage[margin=1in]{geometry}

\newtheorem{hypothesis}{Hypothesis}
\newtheorem{proposition}{Proposition}
\newtheorem{prediction}{Prediction}

\doublespacing

\title{\textbf{Endogenous Stages and Complementary Interventions:\\
A Theory of Behavioral Change Dynamics}}

\author{
[Author Names]\\
Department of Economics, University of Zurich
}

\date{January 2026}

\begin{document}

\maketitle

%%%%%%%%%%%%%%%%%%%%%%%%%%%%%%%%%%%%%%%%%%%%%%%%%%%%%%%%%%%%%%%%%%%%%%%%%%%%%%%
% ABSTRACT
%%%%%%%%%%%%%%%%%%%%%%%%%%%%%%%%%%%%%%%%%%%%%%%%%%%%%%%%%%%%%%%%%%%%%%%%%%%%%%%

\begin{abstract}
\noindent
We develop a formal theory in which behavioral change stages emerge endogenously from awareness, willingness, and action thresholds. Unlike existing stage models that postulate fixed numbers of stages (e.g., five in the Transtheoretical Model), our framework derives stages from primitives and generates testable predictions about stage-dependent intervention effectiveness. The key insight is that intervention mechanisms exhibit \textit{super-additive complementarity}: combining awareness-increasing and threshold-reducing interventions produces effects exceeding their sum. We characterize the conditions under which individuals become ``stuck'' in the contemplation stage ($\kappa_{S3} > 0.7$) and derive the \textit{Habit Lock-In Time} $T_{lock}$---the point after which behavioral change becomes effectively irreversible. The model predicts that stage-targeted interventions outperform uniform interventions by 40--60\%, with the largest gains from matching intervention type to stage position. We propose laboratory and field experiments to test these predictions.

\medskip
\noindent
\textbf{JEL Codes:} D91, D87, C93, I12\\
\textbf{Keywords:} Behavioral change, stage models, intervention design, complementarity, habit formation
\end{abstract}

\newpage

%%%%%%%%%%%%%%%%%%%%%%%%%%%%%%%%%%%%%%%%%%%%%%%%%%%%%%%%%%%%%%%%%%%%%%%%%%%%%%%
% 1. INTRODUCTION
%%%%%%%%%%%%%%%%%%%%%%%%%%%%%%%%%%%%%%%%%%%%%%%%%%%%%%%%%%%%%%%%%%%%%%%%%%%%%%%

\section{Introduction}

Behavioral interventions---from nudges to commitment devices---have become central tools in public policy. Yet a fundamental question remains inadequately addressed: \textit{why do identical interventions produce vastly different effects across individuals?} We propose that the answer lies in stage-dependent dynamics: individuals at different positions in the behavioral change journey respond systematically differently to the same intervention.

Consider two individuals both aware that regular exercise improves health. The first says, ``I know I should exercise, but I haven't really thought about starting.'' The second says, ``I've planned to start running next Monday and already bought shoes.'' Both exhibit positive awareness, but they occupy different \textit{stages} of the behavioral change process. An intervention effective for the second (e.g., a running partner) may entirely fail for the first (who needs to first perceive personal relevance).

This observation motivates our central question: Can we characterize how intervention effectiveness varies systematically with stage position, and can we derive optimal stage-contingent intervention strategies?

\subsection{Contribution}

We make three contributions. First, we develop a formal model in which behavioral change stages emerge \textit{endogenously} from three primitives: awareness $A$, willingness $WAX$, and action threshold $\theta$. Unlike existing models that postulate stages as exogenous constructs, our approach derives them from decision-theoretic foundations.

Second, we characterize intervention mechanisms as operating through distinct channels: awareness-increasing interventions ($\Delta A > 0$), motivation-increasing interventions ($\Delta WAX > 0$), and threshold-reducing interventions ($\Delta \theta < 0$). Critically, we show these mechanisms exhibit \textit{complementarity}: their combined effect exceeds the sum of individual effects.

Third, we derive testable quantitative predictions. We predict the ``contemplation trap'' phenomenon (where $\kappa_{S3} \approx 0.70$ produces stable non-action despite positive intentions), the Habit Lock-In Time $T_{lock}$, and the 40--60\% efficiency gain from stage-targeted versus uniform interventions.

\subsection{Related Literature}

Our work relates to several strands of literature. The Transtheoretical Model \citep{prochaska1983stages} introduced the concept of behavioral change stages but did not derive them from primitives. The behavioral economics literature on nudges \citep{PAP-thaler2008nudge} and commitment devices \citep{bryan2010commitment} has documented intervention effectiveness without systematically addressing stage-dependence.

Most closely related is the literature on complementarity in economic mechanisms. \citet{efferson2024superadditive} demonstrate super-additive cooperation from combining repeated interactions with group competition. \citet{bartling2025complementarity} show that trust and contract enforcement are complements, not substitutes. We extend this logic to behavioral interventions: awareness and friction-reduction are complements in producing behavioral change.


%%%%%%%%%%%%%%%%%%%%%%%%%%%%%%%%%%%%%%%%%%%%%%%%%%%%%%%%%%%%%%%%%%%%%%%%%%%%%%%
% 2. MODEL
%%%%%%%%%%%%%%%%%%%%%%%%%%%%%%%%%%%%%%%%%%%%%%%%%%%%%%%%%%%%%%%%%%%%%%%%%%%%%%%

\section{Model}

\subsection{Primitives}

An agent is characterized by three state variables at time $t$:
\begin{itemize}
    \item \textbf{Awareness} $A(t) \in [0,1]$: The degree to which the behavioral option is mentally available and its consequences understood.
    \item \textbf{Willingness} $WAX(t) \in [0,1]$: The motivational readiness to perform the target behavior.
    \item \textbf{Threshold} $\theta(t) \in [0,\infty)$: The barrier (friction, cost, effort) that must be overcome for action.
\end{itemize}

\begin{proposition}[Action Condition]
Action occurs if and only if $WAX(t) \geq \theta(t)$.
\end{proposition}

The gap $\Delta \equiv WAX - \theta$ determines action probability. When $\Delta < 0$, the agent intends but does not act; when $\Delta > 0$, action is triggered.

\subsection{Endogenous Stage Emergence}

\begin{proposition}[Stage Emergence]
Stages emerge from the joint distribution of $(A, WAX, \theta)$:
\begin{equation}
S(t) = h(A(t), WAX(t), \theta(t), \Psi(t))
\end{equation}
where $\Psi$ is a context vector.
\end{proposition}

We identify six emergent stages:

\begin{table}[h]
\centering
\begin{tabular}{llccc}
\toprule
Stage & Name & $A$ & $WAX$ & $\theta$ \\
\midrule
$S_1$ & Unaware & $\approx 0$ & -- & $\to \infty$ \\
$S_2$ & Aware & low & low & high \\
$S_3$ & Considering & medium & medium & medium \\
$S_4$ & Intending & high & $\approx \theta$ & medium \\
$S_5$ & Acting & high & $> \theta$ & low \\
$S_6$ & Maintaining & -- & -- & $\to 0$ \\
\bottomrule
\end{tabular}
\caption{Emergent stage characteristics}
\label{tab:stages}
\end{table}

The key insight is that stage boundaries are not arbitrary but emerge from discontinuities in the $(A, WAX, \theta)$ space.

\subsection{Transition Dynamics}

Stage transitions follow:
\begin{equation}
\frac{dS}{dt} = \tau \cdot (WAX - \theta) \cdot (1 - \kappa) + \beta \cdot \text{Social} + \varepsilon
\end{equation}

where:
\begin{itemize}
    \item $\tau \in [0.03, 0.25]$ is the transition rate (per month)
    \item $\kappa \in [0.30, 0.85]$ is stage persistence (resistance to leaving current stage)
    \item $\beta \in [0.30, 0.57]$ is social contagion strength
\end{itemize}

\begin{proposition}[Contemplation Trap]
Stage $S_3$ exhibits the highest persistence: $\kappa_{S3} \approx 0.70$. Individuals become ``stuck'' deliberating because the combination of moderate awareness and sub-threshold willingness is locally stable.
\end{proposition}

This explains why many individuals report wanting to change behavior but fail to act: the contemplation stage is an attractor in the dynamics.

\subsection{Relapse}

Backward transitions occur with probability:
\begin{equation}
P(S_i \to S_{i-1}) = \rho \cdot g(\Delta\Psi, t_{since})
\end{equation}

where $\rho \in [0.10, 0.85]$ is the base relapse probability, which varies by domain (highest in health behaviors, lowest in one-time decisions).


%%%%%%%%%%%%%%%%%%%%%%%%%%%%%%%%%%%%%%%%%%%%%%%%%%%%%%%%%%%%%%%%%%%%%%%%%%%%%%%
% 3. INTERVENTION MECHANISMS
%%%%%%%%%%%%%%%%%%%%%%%%%%%%%%%%%%%%%%%%%%%%%%%%%%%%%%%%%%%%%%%%%%%%%%%%%%%%%%%

\section{Intervention Mechanisms and Complementarity}

\subsection{Three Intervention Channels}

Interventions operate through three distinct channels:

\begin{enumerate}
    \item \textbf{Awareness interventions} ($\Delta A > 0$): Information campaigns, salience nudges, social proof (``your neighbors are doing it'')
    \item \textbf{Motivation interventions} ($\Delta WAX > 0$): Incentives, loss framing, goal setting, implementation intentions
    \item \textbf{Friction interventions} ($\Delta \theta < 0$): Defaults, simplification, commitment devices, deadline effects
\end{enumerate}

\subsection{Complementarity}

\begin{hypothesis}[Super-Additive Complementarity]
\label{hyp:complementarity}
The combined effect of awareness and friction interventions exceeds their sum:
\begin{equation}
\Delta B(A + \theta) > \Delta B(A) + \Delta B(\theta)
\end{equation}
where $\Delta B$ denotes behavioral change.
\end{hypothesis}

The intuition is straightforward: awareness without low friction fails because the agent knows but cannot act; low friction without awareness fails because the agent can act but does not know. Together, they unlock behavioral change that neither achieves alone.

This parallels \citet{efferson2024superadditive}'s finding that repeated interaction and group competition exhibit super-additive effects on cooperation.

\subsection{Stage-Dependent Effectiveness}

\begin{hypothesis}[Stage-Intervention Matching]
\label{hyp:matching}
Intervention effectiveness depends on stage:
\begin{itemize}
    \item $S_1 \to S_2$: Awareness interventions most effective
    \item $S_2 \to S_3$: Relevance interventions most effective
    \item $S_3 \to S_4$: Motivation interventions most effective
    \item $S_4 \to S_5$: Friction interventions most effective
    \item $S_5 \to S_6$: Reinforcement interventions most effective
\end{itemize}
\end{hypothesis}

Applying a friction intervention to someone in $S_1$ (Unaware) produces zero effect; the agent needs awareness first. Conversely, applying an awareness intervention to someone in $S_4$ (Intending) is wasteful; they already know but need threshold reduction.


%%%%%%%%%%%%%%%%%%%%%%%%%%%%%%%%%%%%%%%%%%%%%%%%%%%%%%%%%%%%%%%%%%%%%%%%%%%%%%%
% 4. TESTABLE PREDICTIONS
%%%%%%%%%%%%%%%%%%%%%%%%%%%%%%%%%%%%%%%%%%%%%%%%%%%%%%%%%%%%%%%%%%%%%%%%%%%%%%%

\section{Testable Predictions}

\subsection{The Habit Lock-In Time}

\begin{prediction}[Habit Lock-In]
There exists a critical time $T_{lock}$ after which behavioral change becomes effectively irreversible:
\begin{equation}
T_{lock} = \frac{1}{\lambda_D} \cdot \ln\left(\frac{D_0}{\theta_D - z_\varepsilon \cdot \sigma_{RPE,\infty}}\right)
\end{equation}
where $D_0$ is initial doubt, $\lambda_D$ is doubt decay rate, and $\theta_D$ is the redeliberation threshold.
\end{prediction}

For typical parameters, $T_{lock} \approx 66$ days for simple behaviors (consistent with \citealt{lally2010habits}) and 6--12 months for complex behaviors.

\subsection{Efficiency Gains from Targeting}

\begin{prediction}[Targeting Efficiency]
Stage-targeted interventions outperform uniform interventions by 40--60\%:
\begin{equation}
\frac{\text{Effect}_{\text{targeted}}}{\text{Effect}_{\text{uniform}}} \in [1.4, 1.6]
\end{equation}
\end{prediction}

This prediction is directly testable: randomize individuals to receive either (a) a uniform intervention or (b) an intervention matched to their assessed stage.

\subsection{Contemplation Escape}

\begin{prediction}[Breaking the Contemplation Trap]
The combination of social proof ($\Delta A$) and deadline ($\Delta \theta$) increases $S_3 \to S_4$ transition probability by more than either alone:
\begin{equation}
P(S_3 \to S_4 | A + \theta) - P(S_3 \to S_4 | \text{control}) > \sum_i \left[P(S_3 \to S_4 | i) - P(S_3 \to S_4 | \text{control})\right]
\end{equation}
\end{prediction}


%%%%%%%%%%%%%%%%%%%%%%%%%%%%%%%%%%%%%%%%%%%%%%%%%%%%%%%%%%%%%%%%%%%%%%%%%%%%%%%
% 5. EXPERIMENTAL DESIGN
%%%%%%%%%%%%%%%%%%%%%%%%%%%%%%%%%%%%%%%%%%%%%%%%%%%%%%%%%%%%%%%%%%%%%%%%%%%%%%%

\section{Proposed Experimental Tests}

\subsection{Laboratory Experiment: Stage-Intervention Matching}

\textbf{Design:} $6 \times 4$ factorial design.

\textbf{Stage Assessment:} Participants complete a validated questionnaire to determine their current stage for a target behavior (e.g., retirement savings).

\textbf{Intervention Assignment:} Random assignment to one of four conditions:
\begin{enumerate}
    \item Control (no intervention)
    \item Awareness intervention (information provision)
    \item Friction intervention (simplified enrollment)
    \item Combined intervention
\end{enumerate}

\textbf{Outcome:} Behavioral change measured at 1 week, 1 month, and 3 months.

\textbf{Prediction:} Interaction effect between stage and intervention type. Main effects will be positive but smaller than properly matched stage-intervention effects.

\subsection{Field Experiment: Energy Conservation}

\textbf{Setting:} Household energy consumption.

\textbf{Sample:} 2,000 households randomly assigned.

\textbf{Treatment Arms:}
\begin{enumerate}
    \item Control
    \item Uniform nudge (social comparison letter)
    \item Stage-assessed + targeted intervention
    \item Stage-assessed + mismatched intervention (placebo)
\end{enumerate}

\textbf{Prediction:} Treatment 3 (targeted) produces 40--60\% larger effects than Treatment 2 (uniform), and Treatment 4 (mismatched) performs no better than control.

\subsection{Identification Strategy}

Causal identification relies on:
\begin{enumerate}
    \item Random assignment to intervention type
    \item Stage assessment \textit{before} treatment assignment (no selection)
    \item Pre-registered analysis plan
\end{enumerate}

Potential confounds (heterogeneous treatment effects correlated with stage) are addressed by the mismatched-intervention arm, which holds stage assessment constant while varying intervention appropriateness.


%%%%%%%%%%%%%%%%%%%%%%%%%%%%%%%%%%%%%%%%%%%%%%%%%%%%%%%%%%%%%%%%%%%%%%%%%%%%%%%
% 6. DOMAIN-SPECIFIC PREDICTIONS
%%%%%%%%%%%%%%%%%%%%%%%%%%%%%%%%%%%%%%%%%%%%%%%%%%%%%%%%%%%%%%%%%%%%%%%%%%%%%%%

\section{Domain-Specific Predictions}

The model parameters vary by domain. Table \ref{tab:domain} presents calibrated values.

\begin{table}[h]
\centering
\begin{tabular}{lcccc}
\toprule
Domain & $A_{base}$ & $\tau$ & $\rho$ & $\kappa_{S3}$ \\
\midrule
Retirement savings & 0.35 & 0.08/mo & 0.25 & 0.75 \\
Health behaviors & 0.50 & 0.12/mo & 0.45 & 0.70 \\
Energy conservation & 0.30 & 0.18/mo & 0.35 & 0.55 \\
Financial decisions & 0.45 & 0.10/mo & 0.30 & 0.65 \\
\bottomrule
\end{tabular}
\caption{Domain-specific parameters}
\label{tab:domain}
\end{table}

\begin{prediction}[Domain Heterogeneity]
Energy conservation exhibits the fastest transition rates ($\tau = 0.18$) but moderate relapse ($\rho = 0.35$), predicting rapid initial adoption followed by gradual decay without reinforcement. Health behaviors exhibit slower transitions ($\tau = 0.12$) but highest relapse ($\rho = 0.45$), predicting the need for sustained intervention.
\end{prediction}


%%%%%%%%%%%%%%%%%%%%%%%%%%%%%%%%%%%%%%%%%%%%%%%%%%%%%%%%%%%%%%%%%%%%%%%%%%%%%%%
% 7. CONCLUSION
%%%%%%%%%%%%%%%%%%%%%%%%%%%%%%%%%%%%%%%%%%%%%%%%%%%%%%%%%%%%%%%%%%%%%%%%%%%%%%%

\section{Conclusion}

We have developed a theory of behavioral change in which stages emerge endogenously from awareness, willingness, and action thresholds. The framework generates three main predictions: (1) intervention mechanisms exhibit super-additive complementarity, (2) stage-targeted interventions outperform uniform interventions by 40--60\%, and (3) the ``contemplation trap'' is a stable equilibrium that requires combined interventions to escape.

The practical implications are substantial. Current policy approaches typically deploy uniform interventions across heterogeneous populations. Our analysis suggests that substantial efficiency gains---potentially 40--60\%---are available from stage assessment and intervention matching. The cost of stage assessment (a brief questionnaire) is likely far smaller than the efficiency gain from proper targeting.

Future work should test these predictions experimentally and extend the framework to multi-behavior settings where individuals simultaneously pursue multiple behavioral changes with limited cognitive capacity.

%%%%%%%%%%%%%%%%%%%%%%%%%%%%%%%%%%%%%%%%%%%%%%%%%%%%%%%%%%%%%%%%%%%%%%%%%%%%%%%
% REFERENCES
%%%%%%%%%%%%%%%%%%%%%%%%%%%%%%%%%%%%%%%%%%%%%%%%%%%%%%%%%%%%%%%%%%%%%%%%%%%%%%%

\bibliographystyle{aer}

\begin{thebibliography}{99}

\bibitem[Bartling et al.(2025)]{bartling2025complementarity}
Bartling, B., Fehr, E., Huffman, D., \& Netzer, N. (2025).
The complementarity between trust and contract enforcement.
\textit{Economic Journal}, forthcoming.

\bibitem[Bryan et al.(2010)]{bryan2010commitment}
Bryan, G., Karlan, D., \& Nelson, S. (2010).
Commitment devices.
\textit{Annual Review of Economics}, 2, 671--698.

\bibitem[Efferson et al.(2024)]{efferson2024superadditive}
Efferson, C., Bernhard, H., Fischbacher, U., \& Fehr, E. (2024).
Super-additive cooperation.
\textit{Nature}, forthcoming.

\bibitem[Lally et al.(2010)]{lally2010habits}
Lally, P., van Jaarsveld, C. H., Potts, H. W., \& Wardle, J. (2010).
How are habits formed: Modelling habit formation in the real world.
\textit{European Journal of Social Psychology}, 40(6), 998--1009.

\bibitem[Prochaska \& DiClemente(1983)]{prochaska1983stages}
Prochaska, J. O., \& DiClemente, C. C. (1983).
Stages and processes of self-change of smoking.
\textit{Journal of Consulting and Clinical Psychology}, 51(3), 390--395.

\bibitem[Thaler \& Sunstein(2008)]{PAP-thaler2008nudge}
Thaler, R. H., \& Sunstein, C. R. (2008).
\textit{Nudge: Improving decisions about health, wealth, and happiness}.
Yale University Press.

\end{thebibliography}

\end{document}
